\section{Motivation}

Large language model serving is a systems problem as much as a model problem.
An online service must keep expensive accelerators occupied, admit requests
with different prompt and generation lengths, manage a dynamically growing
key--value cache, and coordinate communication across devices. vLLM addresses
several of these concerns through continuous batching and PagedAttention
\cite{kwon2023vllm}. The adoption of sparse \gls{moe} models introduces a
further challenge: the compute and communication performed for a token depend
on which experts the router selects.

An MoE layer contains many feed-forward expert networks but activates only a
small subset for each token. This conditional computation increases model
capacity without making the arithmetic cost proportional to the total number
of parameters \cite{shazeer2017moe,fedus2021switch}. It also makes the runtime
load data dependent. Tokens in a batch do not necessarily select experts
uniformly, and the popularity distribution can shift as the request mix
changes. When experts are sharded across devices with \gls{ep}, an overloaded
expert becomes an overloaded rank. Faster ranks then wait at collective
communication boundaries, so one local imbalance can reduce cluster-wide
throughput.

\gls{eplb} mitigates this problem by observing expert usage and periodically
computing a new physical placement. Popular experts can be replicated and the
mapping between logical experts and physical slots can be changed. Applying a
new placement, however, requires copying expert weights between ranks. These
objects are much larger than control messages and the transfers use the same
host and network resources as foreground inference. This thesis focuses on a
specific but consequential phase of EPLB: scheduling those expert migrations.

The original execution path posts the point-to-point operations for a layer as
one group. This minimizes the number of explicit communication rounds, which
is attractive if migration completion time is the only objective. During
asynchronous EPLB, however, inference does not stop. A rank that simultaneously
sends to and receives from several peers can create a short traffic burst at
its network interface. The burst may compete with token dispatch, result
collection, tensor-parallel communication, or other inference-related traffic.
Thus, the schedule that completes migration soonest is not necessarily the
schedule that gives users the best latency while migration is in progress.

The distinction is especially important in a multi-node deployment whose
network is slower than local accelerator memory. Expert migration is a
background maintenance operation; the foreground objective is serving quality.
The central hypothesis of this work is therefore:

\begin{quote}
Limiting the number of simultaneous migration peers per rank can reduce network
bursts enough to improve foreground serving, even when doing so introduces
additional migration rounds.
\end{quote}

\section{Research questions}

The thesis is organized around four research questions.

\begin{enumerate}[label=\textbf{RQ\arabic*:},leftmargin=2.2cm]
  \item \textbf{How should expert migration contention be represented?}
  A useful model must distinguish expert count, source and destination ranks,
  network paths, and link capacities while remaining simple enough to explain
  the scheduler's behavior.

  \item \textbf{What scheduling rule is suitable for the online path?}
  An exact optimizer is valuable as a design tool, but a production scheduler
  must be deterministic, fast, dependency-light, and compatible with the
  existing replicated placement maps.

  \item \textbf{Can the implementation preserve EPLB semantics?}
  Batching must move every required expert exactly once, retain local-copy and
  replica semantics, respect the communicator lifecycle, and leave the legacy
  path unchanged unless the feature is explicitly enabled.

  \item \textbf{Does reduced migration contention improve serving?}
  The answer requires both user-visible metrics and direct network evidence.
  Throughput, \gls{ttft}, \gls{tpot}, and end-to-end latency are measured
  together with one-second NIC byte counters and scheduler profiling.
\end{enumerate}

\section{Contributions}

This work makes the following contributions.

\begin{itemize}
  \item It separates two objectives that can otherwise be conflated: minimizing
  migration completion time and minimizing the interference experienced by
  concurrent inference.

  \item It formulates contention-aware migration scheduling as an \gls{ilp}.
  The model assigns coalesced directed rank-pair flows to batches, constrains
  each rank to one incident flow per batch, estimates routed link load, and
  minimizes the worst modeled inference time. This is an exact optimum only
  within the hard rank-disjoint-flow policy and the stated linear model. An illustrative
  eight-rank topology exposes the trade-off rather than claiming a calibrated
  latency predictor or proving that rank-disjoint scheduling is universally best.

  \item It designs a production-oriented greedy scheduler. The problem is
  reduced to deterministic first-fit edge coloring of a communication graph.
  Multiple experts with the same directed source--destination pair are
  coalesced, independent rank pairs can progress concurrently, and every rank
  participates in at most one directed flow in a batch.

  \item It integrates the scheduler into asynchronous vLLM EPLB. Scheduling
  for all 48 MoE layers is invoked once per rebalance cycle, while transfers
  remain layer by layer. The feature is opt-in and synchronous EPLB retains the
  original path because inference is already paused in that mode.

  \item It evaluates the complete serving system and profiles the scheduler.
  The final contribution was reviewed and merged into upstream vLLM as pull
  request \prnumber\ \cite{vllmpr52641}.
\end{itemize}

\section{Scope and claims}

The evaluation is deliberately narrow enough to be reproducible. It uses one
four-node cluster, one model, one communication backend, one workload shape,
and one A/B run per condition. This supports a case study and an engineering
validation, not a universal performance claim. The experiment establishes
that the proposed mechanism works and is beneficial under the measured
configuration. It does not establish how the effect scales to modern
datacenter GPUs, RDMA fabrics, different expert sizes, or more ranks.

The thesis also distinguishes artifacts by purpose. The ILP is an offline
analytical baseline. The greedy algorithm is the deployed scheduler. Older
experiments created during design iteration are not used for the final result;
all reported E2E values come from one archived A/B directory, and all scheduler
costs come from a separate archived profiling directory. This provenance is
specified in \cref{ch:methodology} and the appendices.

\section{Thesis structure}

\Cref{ch:background} introduces sparse MoE inference, expert parallelism,
vLLM EPLB, and asynchronous migration. \Cref{ch:related-work} reviews MoE
communication, coflow scheduling, and graph-based migration grouping.
\Cref{ch:formulation,ch:design} then develop the contention model and the
production scheduler within one proposed-system chapter.
\Cref{ch:methodology,ch:results} document the evaluation procedure and report
the E2E, NIC, and profiling results within one experimental-evaluation chapter.
Finally, \cref{ch:conclusions} summarizes the findings, limitations, and future
work.
